% ============================================================
%  第13章 其他高级语言中OpenCL的使用
%  源文件：content/chapter13/
% ============================================================


% ---- 13.1-chinese.md ----
\section{本章简介}

目前为止，我们所讨论的OpenCL所使用的系统编程语言都是C和C++；不过，这只是OpenCL所支持的语言中的一小部分。本章，我么将将来看一看，其他编程语言是如何使用OpenCL进行计算加速的。这些语言有，Java、Python，以及函数式编程语言Haskell。



% ---- 13.2-chinese.md ----
\section{越过C和C++}

对于很多开发者来说，他们会选择C和C++作为开发语言。其他的一些开发者则可能会选择别的编程语言，我们所知的很多软件是使用Java或Python开发的。这些更高级的语言有更高的生产力，通常会提供很多必要的特性(比如，垃圾收集)，而性能并不一定是系统设计者首要考虑的因素。这些语言还有一种特点，可移植性特别好，想想Java的口号“一次编写，处处可用”，这样开发者就不需要为一些系统底层的问题而困扰。但通常很难、有时甚至不可能让这类语言编写的应用达到接近峰值的性能。

为了弥补性能鸿沟，并且不需要将很多已有库使用这类语言重新再写一遍。我们这里会使用\emph{外部函数接口}(FFI，foreignfunctioninterface)，其可以让应用调用使用C和C++(或其他低层语言)编写的原生库。例如，Java提供了Java原生接口，Python也有同样的机制。当然，Java和Python也有已经包装好的OpenCL库(JOCL(Java bindings for OpenCL)[1], PyOpenCL[2])，可以让开发者直接调用。这种模型的确是非常底层，并在托管运行时与OpenCL的原生非托管部分之间提供连接。为了展示这种代码的特点，代码清单13.1中展示了一个PyOpenCL版本的向量相加。

\begin{codeboxt}[python]{代码清单13.1 PyOpenCL实现的向量相加}
import pyopencl as cl
import numpy
import numpy.linalg as la

a = numpy.random.rand(50000).astype(numpy.float32)
b = numpy.random.rand(50000).astype(numpy.float32)

ctx = cl.create_some_context()
queue = cl.CommandQueue(ctx)

mf = cl.mem_flags
a_buf = cl.Buffer(ctx, mf.READ_ONLY | mf.COPY_HOST_PTR, hostbuf = a)
b_buf = cl.Buffer(ctx, mf.READ_ONLY | mf.COPY_HOST_PTR, hostbuf = b)

prg = cl.Program(ctx, """
    __kernel void vecadd(__global const float *a,
                         __global const float *b,
                         __global float *c){
    int gid = get_global_id(0);
    c[gid] = a[gid] + b[gid];
    }
    """).build()

prg.vecadd(queue, a.shape, None, a_buf, b_buf, dest_buf);

a_plus_b = numpy.empty_like(a)
cl.enqueue_copy(queue, a_plus_b, dest_buf)

print la.norm(a_plus_b - (a + b))
\end{codeboxt}

下面的例子使用同样的一种包装API——Aparapi[3]——进行实现(之前是由AMD主导进行开发，不过现在已经是一款开源项目)。Aparapi允许Java开发者使用GPU，将一些可以进行数据并行的代码段放到GPU上运行。Aparapi运行时会把Java字节码转换为OpenCL，再到GPU上执行。如果，因为某些原因Aparapi无法运行在GPU上，其会使用Java线程池进行数据并发的处理。Aparapi的意义在于，保持Java语言的语法，并发其精神。下面我们还会用向量相加展示Aparapi的使用方式(代码清单13.2)，其中没有任何OpenCL C代码，也没有OpenCL API调用。

\begin{codeboxt}[java]{代码清单13.2 使用Aparapi实现的OpenCL，在Java中实现的向量相加}
package com.amd.aparapi.sample.add;
import com.amd.aparapi.Kernel;
import com.amd.aparapi.Range;

public class Main{
  public static void main(String[] _args){
    final int size = 512;
    final float[] a = new float[size];
    final float[] b = new float[size];

    for (int i = 0; i < size; i++){
      a[i] = (float)(Math.random()*100);
      b[i] = (float)(Math.random()*100);
    }

    final float[] sum = new float[size];
    Kenrel kernel = new Kernel(){
      @Override public void run(){
        int gid = getGlobalId();
        sum[gid] = a[gid] + b[gid];
      }
    };

    kernel.execute(Range.create(512));
    for (int i = 0; i < size; i++){
      System.out.printf("%6.2f + %6.2f = %8.2f\n", a[i], b[i], sum[i]);
    }
    kernel.dispose();
  }
}
\end{codeboxt}

另外，Aparapi开发者要是用OpenCL，需要创建出一个Aparapi类的实例，在运行时重载内核函数的实现(内核函数可被动态编译)，并最终由生成的Java字节码动态编译为OpenCL。

Aparapi是把特定领域专属语言(DSL)嵌入到宿主语言中的一个更一般概念的实例，主机端编码语言可以是另外一种语言(本例中主机端是Java)。DSL专注于为该领域的专家提供相应接口，并且DLS通常会为给定的学术领域提供一组特定的特性——例如，医学图像。本例中，使用的领域是通用计算，实现的功能是数据并行计算。

\myRefRule

\begin{myRefs}

[1] Java bindings for OpenCL(JOCL), 2012, \texthref{http://www.jocl.org}

[2] A.Klöckner, \emph{PyOpenCL}, 2012, \texthref{http://mathema.tician.de/software/pyopencl}

[3] Aparapi, 2012, \texthref{http://Aparapi.googlecode.com}

\end{myRefs}



% ---- 13.3-chinese.md ----
\section{Haskell中使用OpenCL}

Haskell是一种纯函数式语言，其属于类\href{https://en.wikipedia.org/wiki/ML\_(programming\_language}{标准ML(SML)})模型函数语言的一种。与其他已经提过的语言不同，Haskell(或SML)编程是通过函数式进行描述，应用会通过表达式的参数对表达式做出对应的判断。总体而言，该模型与命令式编程的区别在于不定义语句的顺序，也不允许副作用。这会使外部声明的值没有进行初始化。这就能看出Haskell类语言的主要优势和劣势。把无副作用的编程与Haskell庞大而复杂的类型系统结合起来，通常会让一些有过C、C++、Java经验的开发者在第一次使用时，有些难以驾驭的感觉。不过，这些问题会在并行程序中解决，例如这个例子，其表达式计算出来的结果是相互独立的，所以其函数式定义是线程安全的。因此，Haskell开发者社区中，逐渐涌现出很多有意思的并行程序。对Haskell感兴趣的读者可以读去一下Hutto写的这本使用Haskell编程的书籍[1]，以及Meijer在微软频道9中的相关教学视频[2]。

通过多方面对Haskell类型系统的论证，不过对于内嵌DSL的设计来说其Haskell类型系统是一个很不错的平台，其能提供相应的抽象化模型，这样的模型能够自动的为GPU编译源码。Accelerate[3]和Obsidian[4]是两个非常不错的例子。不过，本书着重与使用底层OpenCL进行编程，所以我们依旧只会关注，如何让Haskell的编程者通过OpenCL使用GPU。Haskell使用OpenCL，除了带来性能收益，还能获得：

\begin{itemize}
\item OpenCL能为Haskell带来现有CPU线程库无法达到的性能
\item 高级Haskell语言能够减少OpenCL主机端代码的复杂度，可以创建更加高效的开发环境
\end{itemize}

已经有过不止一次为Haskell开发OpenCL包装API的尝试；不过，我们需要使用FFI的方式对OpenCL进行绑定。另外，我们需要更加简单的使用OpenCL，并且还要能发挥OpenCL强大的计算能力。为了达到这一目标，我们会推荐HOpenCL[5]，其为一个开源项目库，提供了对底层OpenCL的包装，并且提供高级的Haskell接口来访问OpenCL API，其消除了很多与OpenCL平台之间的交互，并且比其他的Haskell实现，额外的提供了更强的静态保证。本节剩下的内容，我们将着重与高级的API；不过，对低级API感兴趣的读者，可以去查看HOpenCL的手册和开发者文档。这里需要注意的是，HOpenCL只支持OpenCL 1.2的API。OpenCL 2.0中所添加的新特性，还未加入HOpenCL中。

我们依旧会使用向量相加作为例子。内核代码不会改变，并且直接嵌入到一个字符串中，而需要改变的则是Haskell主机端代码。

\subsection{模块结构}

HOpenCL将实现的一小部分模块放入Langauage.OpenCL结构体中。

\begin{itemize}
\item Language.OpenCL.Host.Constants:定义了OpenCL核心API所使用的基本类型
\item Language.OpenCL.Host.Core:定义了底层OpenCL核心API
\item Language.OpenCL.GLInterop:定义了OpenGL交互的API
\item Language.OpenCL.Host:定义了高级OpenCL API
\end{itemize}

后面的几节中将介绍高级API的相关内容，这里我们将会提到主要API的使用，以及相关的注意事项。对底层实现感兴趣的读者，可以阅读HOpenCL的手册和开发者文档[5]。

\subsection{环境}

很多OpenCL函数需要一个上下文对象，其用来定义OpenCL的执行环境；或是需要一个命令队列对象，提交到队列中的任务将在指定的OpenCL环境中执行。很多OpenCL代码中，有些参数函数作为“噪音”——它们在大部分代码中不会发生变化。为了获得这些信息，HOpenCL提供了两个类，Contextual和Queued。这两个类型的实例可以传入相关的OpenCL API中，执行相应的任务。

通常，使用HOpenCL的应用会使用嵌入计算的方式，这就需要在其他计算式中进行计算——例如，将Queued计算嵌入Contextual计算中，然后尝试将其联系在一起。这里提供的with函数，就是用来完成这件事的：

\begin{codebox}[haskell]
with :: Wraps t m n => t -> m u -> n u
\end{codebox}

\subsection{引用计数}

所有OpenCL对象的声明周期，都不会定义在一个C代码范围内，C API提供对应的操作，手动的减少引用计数(比如：clRetainContext/clReleaseContext)。HOpenCL使用LifeSpan来完成这一概念的定义，并提供相应的retain和release操作：

\begin{codebox}[haskell]
retain :: (LifeSpan t, MonadIO m) => t -> m ()
release :: (LifeSpan t, MonadIO m) => t -> m ()
\end{codebox}

using函数处理构造和释放新(引用计数的)对象。其能够自动对OpenCL对象的声明周期进行管理：

\begin{codebox}[haskell]
using :: (Lifespan t m, CatchIO m) => m t -> (t -> m u) -> m u
\end{codebox}

为了更加简单的使用OpenCL上下文(Context)和命令队列(CommandQueue)，其会自动的在HOpenCL中进行引用计数，withNew操作将with和using的功能融合在一起：

\begin{codebox}[haskell]
withNew :: (Wraps t m n, Lifespan t, CatchIO n) => n t -> m u -> n u
\end{codebox}

\subsection{平台和设备}

与platforms相关的API函数，可以用来在给定系统中，查找可用平台：

\begin{codebox}[haskell]
platforms :: MonadIO m => m [Platform]
\end{codebox}

与C API不同，这里无需对platform查找函数调用两次；HOpenCL将会自动获取全部平台信息。这里唯一麻烦的地方在于，平台信息的返回值为monad m，其是一个MonadIO类的实例，platforms的结果包含在这个实例当中。OpenCL在执行操作时有一定的约束，对于monad对象只能进行输入或输出操作。所有HOpenCL中的OpenCL操作都适用于此限制，这是为了体现底层API可能执行不安全的操作，因此需要顺序执行一些操作。

检查完平台信息之后，可以对(?)操作符进行重载，用来决定我们使用对应平台上的哪种实现。例如，下面的代码就代表我们选择第一个平台作为实现平台，且答应出相应供应商：

\begin{codebox}[haskell]
(p:_) <- platforms
putStrLn .("Platform is by: "++) =<< p ? PlatformVendor
\end{codebox}

通常，任意OpenCL对象信息都需要通过clGetXXXInfo获取，这里的XXX代表着对应OpenCL类型，这里可以这样实现：

\begin{codebox}[haskell]
(?) :: MonadIO m => tv -> qt u -> m u
\end{codebox}

为了平台需要，我们将(?)操作符的类型改一下：

\begin{codebox}[haskell]
(?) :: MonadIO m => Platform -> PlatformInfo u -> m u
\end{codebox}

简单的对OpenCL C++包装API的实现(clGetXXXInfo)，可以通过(?)操作符进行相关信息的返回(需要额外的层提供明确的静态类型)。例如，例子中的PlatformVendor，其返回值的类型就是Haskell中的String类型。

devices函数返回与一个平台相关的一系列设备。其将平台对象和设备类型作为参数传入。设备类型参数可以把设备限定为仅GPU、仅CPU、所有设备，或其它选项。和platforms一样，可以通过(?)操作符对设备信息进行检索：

\begin{codebox}[haskell]
deviceOfType :: MonadIO m => Platform -> [DeviceType] -> m [Device]
\end{codebox}

\subsection{运行环境}

如之前所述，主机需要内核执行在另外一个设备上。为了达到这个目的，上下文对象需要在主机端进行配置，并且需要传入命令和数据到设备端。

\textbf{上下文}

context函数可以根据平台和一组设备对象创建出一个上下文对象：

\begin{codebox}[haskell]
context :: MonadIO m => Platform -> [Device] -> m Context
\end{codebox}

如果需要限定上下文的范围——例如，进行图像交互——然后，通过使用contextFromProperties函数将属性传入上下文：

\begin{codebox}[haskell]
contextFromProperties :: MonadIO m => ContextProperties -> [Device] -> m Context
\end{codebox}

上下文属性也可以传noProperties(其定义了一组空属性值)，pushContextProperty(它向已有的属性集合中添加一个上下文属性)。noProperties和pushContextProperty作为Language.OpenCL.Host.Core结构中的一部分:

\begin{codebox}[haskell]
noProperties :: ContextProperties
pushContextProperty :: ContextProperty t u => t u -> u -> ContextProperties -> ContextProperties
\end{codebox}

\textbf{命令队列}

要向设备提交命令，就需要创建命令队列。queue函数可以通过当前Contextual创建一个命令队列：

\begin{codebox}[haskell]
queuue :: Contextual m => Device -> m CommandQueue
\end{codebox}

由于命令队列是引用计数的，并且定义在特定的Contextual计算中。queue函数的实现通常会合并withNew函数，通过嵌入当前上下文创建命令队列：

\begin{codebox}[haskell]
withNew (queue gpu) $
 __computation dependent on newly created command queue
\end{codebox}

\textbf{内存对象}

buffer函数将会分配一个OpenCL内存对象，并假设其使用的默认标识。函数bufferWithFlags会通过用户指定的内存标识(MemFlag定义在Language.OpenCL.Host.Constatns中)分配一个内存对象：

\begin{codebox}[haskell]
buffer :: (Storable t, Contextural m) => Int -> m (Buffer t)
bufferWithFlags :: (Storable t, Contextual m) => Int -> [MemFlag] -> m (Buffer t)
\end{codebox}

内存对象要和相关的上下文对象相关联，使用using函数可以进行相应的关联操作。

数据从主机传到设备端使用writeTo函数，数据中设备端写回主机端使用readFrom：

\begin{codebox}[haskell]
readFrom :: (Readable cl hs, Storable t, Queued m) => cl t -> Int -> Int -> m (hs t)
writeTo :: (Writable cl hs, Storable t, Queued m) => cl t -> Int -> hs t -> m Event
\end{codebox}

\textbf{创建OpenCL程序对象}

OpenCL程序在运行时可以通过两个函数进行编译，programFromSource和buildProgram。先通过源码创建一个OpenCL程序对象，然后对程序对象进行编译。

\begin{codebox}[haskell]
programFromSource :: Contextual m => String -> m Program
buildProgram :: MonadIO m => Program -> [Device] -> String -> m()
\end{codebox}

\textbf{OpenCL内核}

内核通过函数kenrel创建：

\begin{codebox}[haskell]
kernel :: MonadIO m => Program -> String -> m Kernel
\end{codebox}

参数需要逐个通过函数fixArgument传入。不过，通常参数会在内核在调用前在进行参数传递，并且HOpenCL提供内核invoke函数：

\begin{codebox}[haskell]
fixArgument :: (KernelArgument a, MonadIO m) => Kernel -> Int -> a -> m()
invoke :: KernelInvocation r => kernel -> r
\end{codebox}

HOpenCL还提供了另外一种内核调用方式，其可以将内核认为是闭合的，通过setArgs函数对内核的参数进行设置(这种方式在多线程的上下文中十分有用)：

\begin{codebox}[haskell]
setArgs :: Kernel -> [Co.kernel -> Int -> IO ()] -> Invocation
\end{codebox}

通过一次调用invoke函数，并不能能够完全将一个内核入队；因此，invoke函数需要和overRange函数一起使用，其会将执行域和结果作为一个事件进行入队：

\begin{codebox}[haskell]
overRange :: Queued m => Invocation -> ([Int], [Int], [Int]) -> m Event
\end{codebox}

\textbf{向量相加的实现源码}

下面就是使用HOpenCL实现的向量相加源码：

\begin{codebox}[haskell]
module VecAdd where

import Language.OpenCL.Host
import Language.OpenCL.Host.FFI

import Control.Monad.Trans (lift0)

source =
"__kernel void vec add                                 \n" ++
"  __global int *C, __global int *A, __global int *B){ \n" ++
"  int tid = get_global_id(0);                         \n" ++
" C[tid] = A[tid] + B[tid];                            \n" ++
"}"

elements = 2048 :: Int

main = do (p:_) <- platforms
          [gpu] <- devicesOfType p [GPU]
          withNew (context p [gpu]) $
              using (programFromSource source) $ \p ->
              using (buffer elements) $ \inBufA ->
              using (buffer elements) $ \inBufB ->
              using (buffer elements) $ \outBuf ->
                    do { buildProgram p [gpu] ""
                       ; using (kenrel p "vecadd") $ \vecadd ->
                         withNew (queue gpu) $
                           do writeTo inBufA 0 [0.. elements - 1]
                              writeTo inBufB 0 [0.. elements - 1]
                              invoke vecadd outBuf inBufA inBufB
                                    'overRange' ([0], [elements], [1])
                              (x::[Int]) <- readFrom outBuf 0 elements
                              liftIO (if and $ zipWith (\a b -> a == b+b))
                                                    x [0.. elements - 1]
                                  then print "Output is correct"
                                  else print "Output is incorrect")
                    }
\end{codebox}

\myRefRule

\begin{myRefs}

[1] G. Hutton. Programming in Haskell, Cambridge University Press, Cambridge, 2007.

[2] E. Meijer, Functional Programming Fundamentals. Channel 9 Lectures, 2009. \texthref{http://channel9.msdn.com/Series/C9-Lectures-Erik-Meijer-Functional-Programming-Fundamentals/Lecture-Series-Erik-Meijer-Functional-Programming-Fundamentals-Chapter-1}

[3] M.M. Chakravarty, G. Keller, S. Lee, T.L. McDonell, V.Grover, Accelerating Haskell array codes with multicore GPUs, in: Processdings on the Sixth Workshop on Declarative Aspects of Multicore Programming, ACM DAMP'11, New York, NY, 2011, pp.3-14

[4] J. Svensson, M. Sheeran, K. Claessen, Obsidian: a domain specific embedded language for parallel programming of graphics processors. in: S.-B. Scholz, O. Chitil(Eds), Implementation and Application of Functional Languages, Lecture Notes in Computer Science, vol.5836, Springer, Berlin/Heidelberg, 2011, pp.156-173

[5] B.R Gaster, J. Garrett Morris, HOpenCL, 2012, \texthref{https://github.com/bgaster/hopencl.git}

\end{myRefs}



% ---- 13.4-chinese.md ----
\section{本章总结}

本章，我们看到了在不使用C和C++编程语言，使用其他编程语言的情况下，我们依旧可以使用OpenCL。我们指出许多语言(包括Java和Python)都有可用于生产的OpenCL绑定。还关注了一下函数编程语言是如何进行OpenCL编程的，并用Haskell完成了一个例子。


